% block: shifted-complex-symmetric-heading page:183 type:section
\subsection{{移位复对称线性方程组}}\label{sec:shifted-complex-symmetric-systems}
{本节始终假设 \eqref{eq:shifted-linear-systems} 中的系数矩阵 $A^{(k)}$ 为复对称矩阵, 即 $A^{(k)}=(A^{(k)})^{\top}\ne(A^{(k)})^{\mathrm H}$.}
% block: shifted-cocg-heading page:183 type:subsection
\subsubsection{{移位 COCG 方法}}\label{subsec:shifted-cocg}
{移位 COCG 方法 \BookCite{184} 的推导与 \eqref{eq:shifted-cg-residual-collinearity}--\eqref{eq:shifted-cg-shared-direction} 相同. 移位 CG 方法与移位 COCG 方法的唯一区别是 $\alpha_n$ 和 $\beta_n$. 移位 COCG 方法的算法见 Algorithm\ \ref{alg:shifted-cocg}.}
% block: shifted-cocg-algorithm page:183 type:ordinary-paragraph
\begin{algorithm}[H]\caption{{移位 COCG 方法 (注意: $(\bar{\bm a},\bm b)=\bm a^{\top}\bm b$)}}\label{alg:shifted-cocg}\begin{algorithmic}[1]
\Statex\textbf{Input:} {选择种子系统 $s\in S=\{1,2,\ldots,m\}$ 并令 $\bm r_0^{(s)}=\bm b,\ \beta_{-1}^{(s)}=0$}
\Statex\textbf{Input:} $\bm x_0^{(k)}=\bm p_{-1}^{(k)}=\bm0,\ \pi_0^{(s,k)}=\pi_{-1}^{(s,k)}=1$ {对于} $k=1,2,\ldots,m$
\Statex\textbf{Output:} $\bm x_n^{(k)}$ {对于} $k=1,2,\ldots,m$
\For{$n=0,1,\ldots$, {直至收敛}}
\State $\bm p_n^{(s)}=\bm r_n^{(s)}+\beta_{n-1}^{(s)}\bm p_{n-1}^{(s)}$
\State $\alpha_n^{(s)}=\dfrac{(\bar{\bm r}_n^{(s)},\bm r_n^{(s)})}{(\bar{\bm p}_n^{(s)},A^{(s)}\bm p_n^{(s)})}$
\State $\bm x_{n+1}^{(s)}=\bm x_n^{(s)}+\alpha_n^{(s)}\bm p_n^{(s)}$
\State \{{移位系统开始}\}
\State {执行移位 CG 方法 (Algorithm\ \ref{alg:shifted-cg}) 的第 6--15 行.}
\State \{{移位系统结束}\}
\State $\bm r_{n+1}^{(s)}=\bm r_n^{(s)}-\alpha_n^{(s)}A^{(s)}\bm p_n^{(s)}$
\State $\beta_n^{(s)}=\dfrac{(\bar{\bm r}_{n+1}^{(s)},\bm r_{n+1}^{(s)})}{(\bar{\bm r}_n^{(s)},\bm r_n^{(s)})}$
\EndFor\end{algorithmic}\end{algorithm}
% block: shifted-cocg-properties pages:183,184 type:ordinary-paragraph
{Algorithm\ \ref{alg:shifted-cocg} 具有下列性质:}
\begin{enumerate}
\item {需要计算 $A^{(s)}$ 与向量的乘积, 但不需要计算其他 $k$ 对应的 $A^{(k)}$ 与向量的乘积;}
\item {移位 COCG 方法的近似解 $\bm x_n^{(k)}$ 与将 COCG 方法 (以 $\bm x_0=\bm0$) 应用于 $A^{(k)}\bm x^{(k)}=\bm b$ 得到的第 $n$ 个近似解相同.}
\end{enumerate}
{当对所有 $k$ 矩阵 $A^{(k)}$ 都为实对称矩阵时, 移位 COCG 方法等价于移位 CG 方法.}
% block: generalized-shifted-cocg page:184 type:ordinary-paragraph
{文献 \BookCite{190} 将移位 COCG 方法推广到求解下列线性方程组:}
\begin{equation}\label{eq:generalized-shifted-systems}(A+\sigma_kB)\bm x^{(k)}=\bm b,\end{equation}
{它们称为广义移位线性方程组. 这里 $A$ 和 $B$ 为实对称矩阵, $\sigma_k$ 为复数, 因而系数矩阵 $A+\sigma_kB$ 为复对称矩阵. 容易看出, 若 $B$ 为单位矩阵, 则广义移位线性方程组 \eqref{eq:generalized-shifted-systems} 退化为移位线性方程组.}
{注意, 移位 COCG 方法和其他移位 Krylov 子空间方法不能应用于 \eqref{eq:generalized-shifted-systems}, 因为移位不变性不再成立, 即当 $i\ne j$ 时 $\mathcal K_n(A+\sigma_iB,\bm b)\ne\mathcal K_n(A+\sigma_jB,\bm b)$. 一种补救方法是考虑下列移位线性方程组:}
\[(B^{-1}A+\sigma_k I)\bm x^{(k)}=B^{-1}\bm b,\]
{其中假设 $B$ 非奇异. 然而, $B^{-1}A+\sigma_k I$ 为非 Hermitian 矩阵, 因而可以使用非 Hermitian 线性方程组的移位 Krylov 子空间方法. 另一方面, 已有研究表明, 无需采用这种通用求解器, 就能构造与移位 COCG 方法类似的算法, 称为广义移位 COCG 方法 \BookCite{190}. 若 $B=I$, 则广义移位 COCG 方法退化为移位 COCG 方法.}
% block: generalized-shifted-applications page:184 type:ordinary-paragraph
{如 Section\ \ref{subsec:electronic-structure} 所指出, 若 Hamiltonian 矩阵为实对称矩阵, 则移位线性方程组的系数矩阵 $(\epsilon+i\delta)I-H$ 为复对称矩阵. 因此, 移位 COCG 方法是一种可选方法, 其在计算物理中的应用例如见 \BookCite{73}, \BookCite{110}, \BookCite{111}. 文献 \BookCite{100} 开发了移位 COCG 方法及相关求解器的开源库.}
{广义移位线性方程组也出现在下列线性方程组中:}
\begin{equation}\label{eq:quadratic-parameter-linear-system}(\sigma^2A+\sigma B+C)\bm x=\bm b,\end{equation}
{其中 $A,B,C$ 为复对称矩阵. 文献 \BookCite{158}, \BookCite{159} 考虑了这种类型的参数化线性方程组及其在结构动力学中的应用. 文献 \BookCite{158}, \S3 中一个有趣的观察是, 可以将 \eqref{eq:quadratic-parameter-linear-system} 变换为下列线性方程组:}
\begin{equation}\label{eq:quadratic-system-shifted-linearization}\left(\begin{bmatrix}B&C\\C^{\top}&O\end{bmatrix}+\sigma\begin{bmatrix}A&O\\O&-C^{\top}\end{bmatrix}\right)\begin{bmatrix}\bm y\\\bm x\end{bmatrix}=\begin{bmatrix}\bm b\\\bm0\end{bmatrix}\end{equation}
{前提是 $C$ 非奇异. 由于 $A,B,C$ 为复对称矩阵, \eqref{eq:quadratic-system-shifted-linearization} 的系数矩阵构成复对称的广义移位线性方程组. 因此, 广义移位 COCG 方法也可用于求解 \eqref{eq:quadratic-parameter-linear-system}.}
% block: seed-switching-heading page:185 type:subsection
\paragraph{{种子切换技术}}\label{subsubsec:seed-switching}\mbox{}\par
{由 Algorithm\ \ref{alg:shifted-cocg} (另见 Algorithm\ \ref{alg:shifted-cg} 的第 8 行) 可知, 若 $|\pi_n^{(s,k)}|=|R_n^{(s)}(\sigma_s-\sigma_k)|\ge1$, 则 $\|\bm r_n^{(k)}\|\le\|\bm r_n^{(s)}\|$. 因而, 若能找到满足 $|R_n^{(s)}(\sigma_s-\sigma_k)|\ge1$ 的种子系统, 就可求解所有移位系统. 然而, 除 \BookCite{67}, Corollary 1 中讨论的一些特殊情形外, 预先找到这样的系统极其困难. 种子切换技术 \BookCite{166} 可以避免这一问题, 包含下列步骤:}
\begin{enumerate}[label=(S\arabic*)]
\item {选择一个种子系统, 然后启动 Algorithm\ \ref{alg:shifted-cocg}.}
\item {若在第 $n$ 次迭代已求解种子系统, 则从尚未求解的系统中选取一个新种子系统.}
\item {利用新种子系统从第 $(n+1)$ 次迭代开始执行 Algorithm\ \ref{alg:shifted-cocg}.}
\end{enumerate}
% block: seed-switching-first-stage page:185 type:ordinary-paragraph
{在 (S2) 中, 选取新种子系统 $\widetilde s$ 的一个准则可以是}
\[\widetilde s=\operatorname*{arg\,max}_{i\in I}\{\|\bm r_n^{(i)}\|\},\]
{其中 $I$ 表示尚未求解的线性方程组的指标集. 在 (S3) 中, 需要两个步骤将旧种子系统切换为新种子系统. 首先, 计算}
\[\pi_{n+1}^{(s,\widetilde s)}=R_{n+1}^{(s)}(\sigma_s-\sigma_{\widetilde s}),\quad\beta_n^{(\widetilde s)}=\left(\pi_n^{(s,\widetilde s)}/\pi_{n+1}^{(s,\widetilde s)}\right)^2\beta_n\]
{以得到 $\bm r_{n+1}^{(\widetilde s)}$ 和 $\beta_n^{(\widetilde s)}\bm p_n^{(\widetilde s)}$. 由于从 $\bm r_{n+1}^{(\widetilde s)}+\beta_n^{(\widetilde s)}\bm p_n^{(\widetilde s)}$ 可以得到 $\bm p_{n+1}^{(\widetilde s)}$, 因而可以从第 $(n+1)$ 次迭代开始, 用 COCG 方法求解系统 $(A+\sigma_{\widetilde s}I)\bm x^{(\widetilde s)}=\bm b$. 其次, 为利用新种子 $\widetilde s$ 求解其余线性方程组, 必须从新种子生成参数 $\alpha_{n+1}^{(i)}$ 和 $\beta_n^{(i)}$. 可见, 利用下列多项式很容易生成这些参数:}
\[\pi_{n+1}^{(\widetilde s,i)}=R_{n+1}^{(\widetilde s)}(\sigma_{\widetilde s}-\sigma_i)\quad\text{for all }i\in I.\]
% block: seed-switching-second-stage page:185 type:ordinary-paragraph
{为得到上述多项式, 需要计算}
\[\alpha_i^{(\widetilde s)}=\left(\pi_i^{(s,\widetilde s)}/\pi_{i+1}^{(s,\widetilde s)}\right)\alpha_i,\quad\beta_j^{(\widetilde s)}=\left(\pi_j^{(s,\widetilde s)}/\pi_{j+1}^{(s,\widetilde s)}\right)^2\beta_j\]
{其中 $i=0,\ldots,n$, $j=0,\ldots,n-1$. 因此, 切换策略只需标量运算, 而且可见, 若不发生中断, 反复执行 (S2) 至 (S3) 的过程, 就能在不损失上一次切换之前已生成的 Krylov 子空间维数的情况下继续求解这些系统. 种子切换技术使我们摆脱种子系统选择问题, 也可应用于 Section\ \ref{subsec:shifted-cg} 中的移位 CG 方法和 Section\ \ref{subsec:shifted-bicg} 中的移位 BiCG 方法, 因为移位系统的移位 CG (BiCG) 方法的残差向量就是真正的 CG (BiCG) 残差.}
% block: shifted-cocr-heading page:186 type:subsection
\subsubsection{{移位 COCR 方法}}\label{subsec:shifted-cocr}
{移位 COCR 方法 \BookCite{172} 的推导也与 \eqref{eq:shifted-cg-residual-collinearity}--\eqref{eq:shifted-cg-shared-direction} 相同. 移位 CR 方法与移位 COCR 方法的唯一区别是 $\alpha_n$ 和 $\beta_n$. 移位 COCR 方法的算法见 Algorithm\ \ref{alg:shifted-cocr}.}
% block: shifted-cocr-properties page:186 type:ordinary-paragraph
{Algorithm\ \ref{alg:shifted-cocr} 具有下列性质:}
\begin{enumerate}
\item {需要计算 $A^{(s)}$ 与向量的乘积, 但不需要计算其他 $k$ 对应的 $A^{(k)}$ 与向量的乘积, 这一性质与移位 COCG 方法相同;}
\item {与移位 COCG 方法不同, 移位 COCR 方法的 $\bm x_n^{(k)}$ 与将 COCR 方法 (以 $\bm x_0=\bm0$) 应用于 $A^{(k)}\bm x^{(k)}=\bm b$ 得到的第 $n$ 个近似解不同.}
\end{enumerate}
% block: shifted-cocr-algorithm page:186 type:ordinary-paragraph
\begin{algorithm}[H]\caption{{移位 COCR 方法 (注意: $(\bar{\bm a},\bm b)=\bm a^{\top}\bm b$)}}\label{alg:shifted-cocr}\begin{algorithmic}[1]
\Statex\textbf{Input:} {选择种子系统 $s\in S=\{1,2,\ldots,m\}$ 并令 $\bm r_0^{(s)}=\bm b,\ \beta_{-1}^{(s)}=0$}
\Statex\textbf{Input:} $\bm x_0^{(k)}=\bm p_{-1}^{(k)}=\bm0,\ \pi_0^{(s,k)}=\pi_{-1}^{(s,k)}=1$ {对于} $k=1,2,\ldots,m$
\Statex\textbf{Output:} $\bm x_n^{(k)}$ {对于} $k=1,2,\ldots,m$
\For{$n=0,1,\ldots$, {直至收敛}}
\State $\bm p_n^{(s)}=\bm r_n^{(s)}+\beta_{n-1}^{(s)}\bm p_{n-1}^{(s)}$
\State $(A^{(s)}\bm p_n^{(s)}=A^{(s)}\bm r_n^{(s)}+\beta_{n-1}^{(s)}A^{(s)}\bm p_{n-1}^{(s)})$
\State $\alpha_n^{(s)}=\dfrac{(\bar{\bm r}_n^{(s)},A^{(s)}\bm r_n^{(s)})}{(\overline{A^{(\sigma_s)}\bm p_n^{(s)}},A^{(s)}\bm p_n^{(s)})}$
\State $\bm x_{n+1}^{(s)}=\bm x_n^{(s)}+\alpha_n^{(s)}\bm p_n^{(s)}$
\State \{{移位系统开始}\}
\State {执行移位 CG 方法 (Algorithm\ \ref{alg:shifted-cg}) 的第 6--15 行.}
\State \{{移位系统结束}\}
\State $\bm r_{n+1}^{(s)}=\bm r_n^{(s)}-\alpha_n^{(s)}A^{(s)}\bm p_n^{(s)}$
\State $\beta_n^{(s)}=\dfrac{(\bar{\bm r}_{n+1}^{(s)},A^{(s)}\bm r_{n+1}^{(s)})}{(\bar{\bm r}_n^{(s)},A^{(s)}\bm r_n^{(s)})}$
\EndFor\end{algorithmic}\end{algorithm}
% block: shifted-qmr-sym-heading page:186 type:subsection
{当对所有 $k$ 矩阵 $A^{(k)}$ 都为实对称矩阵时, 移位 COCR 方法等价于移位 CR 方法.}
\subsubsection{{移位 QMR\_SYM 方法}}\label{subsec:shifted-qmr-sym}
{文献 \BookCite{62} 提出了求解非 Hermitian 线性方程组的移位 QMR 方法. 如 \BookCite{61} 所知, 移位 QMR\_SYM 方法是移位 QMR 方法的简化, 且移位 QMR\_SYM 的推导与移位 MINRES 方法类似.}
% block: shifted-qmr-sym-projection page:187 type:ordinary-paragraph
{由 Section\ \ref{subsec:complex-symmetric-lanczos} 中的复对称 Lanczos 过程, $\bm x_n\in\mathcal K_n(A,\bm b)$ 可改写为}
\begin{equation}\label{eq:shifted-qmr-sym-coordinate-solution}\bm x_n^{(k)}=V_n\bm y_n^{(k)},\quad k=1,2,\ldots,m,\end{equation}
{其中 $\bm y_n^{(k)}\in\mathbb C^n$. 对于 \eqref{eq:shifted-linear-systems}, 相应残差向量为 $\bm r_n^{(k)}:=\bm b-(A+\sigma_k I)\bm x_n^{(k)}$. 由 \eqref{eq:shifted-qmr-sym-coordinate-solution} 和复对称 Lanczos 过程的矩阵形式 \eqref{eq:complex-lanczos-matrix-relation}, 得}
\begin{equation}\label{eq:shifted-qmr-sym-projected-residual}\begin{aligned}\bm r_n^{(k)}&=\bm b-(A+\sigma_k I)V_n\bm y_n^{(k)}\\&=V_{n+1}g_1\bm e_1-(AV_n+\sigma_kV_n)\bm y_n^{(k)}\\&=V_{n+1}(g_1\bm e_1-T_{n+1,n}^{(k)}\bm y_n^{(k)}),\quad T_{n+1,n}^{(k)}:=T_{n+1,n}+\sigma_k\begin{bmatrix}I_n\\\bm0^{\top}\end{bmatrix}.\end{aligned}\end{equation}
{这里 $\bm e_1=(1,0,\ldots,0)^{\top}$ 为第一个单位向量, $g_1=(\bar{\bm b},\bm b)^{1/2}$. 与 Section\ \ref{subsec:qmr-sym} 中的 QMR\_SYM 方法一样, $\bm y_n^{(k)}$ 通过求解下列极小化问题确定.}
\begin{equation}\label{eq:shifted-qmr-sym-least-squares}\min_{\bm y_n^{(k)}\in\mathbb C^n}\|g_1\bm e_1-T_{n+1,n}^{(k)}\bm y_n^{(k)}\|.\end{equation}
{极小化问题 \eqref{eq:shifted-qmr-sym-least-squares} 可利用 Section\ \ref{subsec:minres} 中介绍的 Givens 旋转求解. 移位 QMR\_SYM 方法的算法见 Algorithm\ \ref{alg:shifted-qmr-sym}.}
% block: shifted-qmr-sym-properties page:187 type:ordinary-paragraph
{Algorithm\ \ref{alg:shifted-qmr-sym} 具有下列性质:}
\begin{enumerate}
\item {需要计算 $A^{(s)}$ 与向量的乘积, 但不需要计算其他 $k$ 对应的 $A^{(k)}$ 与向量的乘积, 这一性质与移位 CG 方法和移位 CR 方法相同;}
\item {移位 QMR\_SYM 方法的近似解 $\bm x_n^{(k)}$ 与将 QMR\_SYM 方法 (以 $\bm x_0=\bm0$) 应用于 $A^{(k)}\bm x^{(k)}=\bm b$ 得到的第 $n$ 个近似解相同.}
\end{enumerate}
% block: shifted-qmr-sym-algorithm page:187 type:ordinary-paragraph
\begin{algorithm}[H]\caption{{移位 QMR\_SYM 方法 (注意: $(\bar{\bm a},\bm b)=\bm a^{\top}\bm b$)}}\label{alg:shifted-qmr-sym}\begin{algorithmic}[1]
\Statex\textbf{Input:} $\bm x_0^{(k)}=\bm p_{-1}^{(k)}=\bm p_0^{(k)}=\bm0,\ \bm v_1=\bm b/(\bar{\bm b},\bm b)^{1/2},\ g_1^{(k)}=(\bar{\bm b},\bm b)^{1/2},\ \beta_{-1}=0$
\Statex\textbf{Output:} $\bm x_n^{(k)}$ {对于} $l=1,2,\ldots,m$
\For{$n=1,2,\ldots$}
\State ({复对称 Lanczos 过程})
\State $\alpha_n=(\bar{\bm v}_n,A\bm v_n)$
\State $\widetilde{\bm v}_{n+1}=A\bm v_n-\alpha_n\bm v_n-\beta_{n-1}\bm v_{n-1}$
\State $\beta_n=(\overline{\widetilde{\bm v}_{n+1}},\widetilde{\bm v}_{n+1})^{1/2}$
\State $\bm v_{n+1}=\widetilde{\bm v}_{n+1}/\beta_n$
\State $t_{n-1,n}^{(\ell)}=\beta_{n-1},\ t_{n,n}^{(\ell)}=\alpha_n+\sigma_{\ell},\ t_{n+1,n}^{(\ell)}=\beta_n$
\State ({用 Givens 旋转求解最小二乘问题})
\State {执行移位 MINRES 方法 (Algorithm\ \ref{alg:shifted-minres}) 的第 9--24 行.}
\EndFor\end{algorithmic}\end{algorithm}
% block: hermitian-complex-shift-qmr page:188 type:ordinary-paragraph
{文献 \BookCite{60} 考虑了下列移位线性方程组:}
\begin{equation}\label{eq:hermitian-matrix-complex-shift}(A+\sigma_k I)\bm x^{(k)}=\bm b,\end{equation}
{其中 $A$ 为 Hermitian 矩阵, $\sigma_k$ 为复数. 因而系数矩阵 $(A+\sigma_k I)$ 既非 Hermitian, 也非复对称. 在这种情况下, 可以使用非 Hermitian 移位线性方程组的移位 QMR 方法. 另一方面, 文献 \BookCite{60} 中有一个比移位 QMR 方法更高效的算法, 它不需要计算 $(A+\sigma_k I)^{\mathrm H}$ 与向量的乘积, 即每步迭代所需的矩阵向量乘法次数仅为一次. 其关键思想是用 Lanczos 过程生成 Krylov 子空间 $\mathcal K_n(A,\bm b)$, 并利用其基向量求解 \eqref{eq:hermitian-matrix-complex-shift}. 注意, 若不生成 $\mathcal K_n(A,\bm b)$, 而生成 $\sigma$ 为复数时的 $\mathcal K_n(A+\sigma I,\bm b)$, 则需要双 Lanczos 过程; 在矩阵向量乘法最耗时的情况下, 其计算开销约为 Lanczos 过程的两倍.}
% block: generalized-shifted-qmr page:188 type:ordinary-paragraph
{文献 \BookCite{168} 将移位 QMR\_SYM 方法推广到求解 (复对称) 广义移位线性方程组 \eqref{eq:generalized-shifted-systems}. 与 Section\ \ref{subsec:shifted-cocg} 所述的广义移位 COCG 方法类似, 当 $B=I$ 时, 广义移位 QMR\_SYM 方法 \BookCite{168} 退化为移位 QMR\_SYM 方法. 此外, 文献 \BookCite{168} 基于 \BookCite{167} 中用于复对称移位线性方程组的移位加权 QMR\_SYM 方法, 提出了求解复对称广义移位线性方程组的相应算法.}
